\chapter{Testosterone}

\section*{What Is This Test?}
Testosterone is a C19 steroid hormone and the primary androgen in males. It is produced primarily by the Leydig cells of the testes (95\% in males) under LH stimulation, with minor contributions from the adrenal cortex (zona reticularis) and peripheral conversion of androstenedione. In females, testosterone is produced by the ovarian theca cells and adrenal cortex (~50\% each), plus peripheral conversion of androstenedione. Testosterone circulates in blood in three forms: (1) tightly bound to sex hormone-binding globulin (SHBG, ~44--60\%), (2) weakly bound to albumin (~38--54\%), and (3) free (unbound, ~1--2\%). The free testosterone plus albumin-bound testosterone constitute bioavailable testosterone, which can dissociate and enter target cells. SHBG-bound testosterone is tightly bound and essentially inactive. Testosterone acts either directly on the androgen receptor or after 5-alpha-reduction to dihydrotestosterone (DHT, a more potent androgen) or aromatization to estradiol. Testosterone is essential for: male sexual differentiation in utero, development of secondary sexual characteristics at puberty, maintenance of libido and erectile function, spermatogenesis, muscle mass and strength, bone density, erythropoiesis (stimulates erythropoietin), and mood/energy. The testosterone test may measure total testosterone (bound plus free), free testosterone (by equilibrium dialysis, the gold standard, or calculated from total testosterone, SHBG, and albumin), or bioavailable testosterone. Total testosterone is the initial screening test for hypogonadism in men and hyperandrogenism in women.

\section*{Alternative Names}
\begin{itemize}
  \item Total Testosterone
  \item Free Testosterone (Equilibrium Dialysis or Calculated)
  \item Bioavailable Testosterone
  \item Serum Testosterone
  \item Plasma Testosterone
  \item Testosterone, Total and Free
  \item Androgen Level
  \item Testosterone by LC-MS/MS
\end{itemize}
\section*{Principle}
Total testosterone is measured using competitive immunoassays on automated analyzers. An anti-testosterone antibody competes with testosterone in the sample for labeled testosterone (chemiluminescent tracer). A displacement agent releases testosterone from SHBG and albumin. The signal is inversely proportional to testosterone concentration. Common platforms: Roche Cobas ECLIA, Abbott Architect CMIA, Siemens Atellica CLIA. Testosterone immunoassays have significant limitations: poor accuracy at low testosterone concentrations (<150 ng/dL in men, and most values in women and children), cross-reactivity with DHT (2--5\%) and other androgens, and inter-assay variability (CV 10--20\% at low concentrations). The gold standard reference method is liquid chromatography-tandem mass spectrometry (LC-MS/MS), which offers superior accuracy, sensitivity (<1 ng/dL), and specificity by separating testosterone from interfering steroids. The 2010 Endocrine Society guidelines on testosterone therapy and the 2013 CDC Hormone Standardization Program recommend LC-MS/MS as the preferred method for testosterone measurement, especially in women, children, and hypogonadal men where immunoassays are inaccurate. Free testosterone is best measured by equilibrium dialysis (gold standard), where free testosterone diffuses across a semipermeable membrane and is measured by LC-MS/MS or RIA. Calculated free testosterone uses total testosterone, SHBG, and albumin in mass action equations (Vermeulen equation) and correlates well with equilibrium dialysis. Direct free testosterone analog immunoassays are inaccurate and NOT recommended.

\section*{History}
Testosterone was first isolated from bull testes by Fred Koch and colleagues at the University of Chicago in 1927. Its chemical structure was determined by Adolf Butenandt in 1935 (Nobel Prize 1939). The first therapeutic use of testosterone (as testosterone propionate injection) was reported by Kenyon in 1937 for hypogonadal men. The first radioimmunoassay for testosterone was developed by Guy Abraham in 1969. The role of 5-alpha-reductase in converting testosterone to DHT was discovered in the 1960s. The concept of measuring free and bioavailable testosterone and the importance of SHBG was elucidated in the 1970s--1980s by Rosner, Vermeulen, and others. LC-MS/MS methods for testosterone were developed in the 1990s and became the recommended standard in the 2010s. The 2010 Endocrine Society clinical practice guideline on testosterone therapy in androgen deficiency syndromes (updated 2018) standardized the diagnostic approach and treatment of male hypogonadism. The 2014 Endocrine Society guideline on androgen therapy in women addressed the controversial role of testosterone in female sexual dysfunction.

\section*{Why This Test Is Ordered}
\begin{itemize}
  \item Diagnosis of male hypogonadism in men with compatible symptoms or signs and consistently low testosterone --- a commonly used threshold is total testosterone <300 ng/dL, but diagnosis requires two separate early-morning measurements and laboratory-specific reference ranges; symptoms and biochemical results must be interpreted together \cite{bhasin2018testosterone}
  \item Classification of hypogonadism as primary (low testosterone, elevated LH/FSH) vs. secondary (low testosterone, inappropriately normal or low LH/FSH)
  \item Monitoring testosterone replacement therapy in hypogonadal men --- aim for a formulation-specific mid-normal range, with sampling time determined by the formulation and prescribing guideline; do not interpret a result without documenting the dose and timing \cite{bhasin2018testosterone}
  \item Evaluation of hyperandrogenism in women with hirsutism, severe inflammatory acne, androgenic alopecia, oligomenorrhea/amenorrhea, or virilization (clitoromegaly, deepening voice, temporal hair recession) --- elevated total or free testosterone supports PCOS or androgen-secreting tumor
  \item Diagnosis of polycystic ovary syndrome (PCOS) --- elevated total or free testosterone is a key diagnostic criterion in the Rotterdam criteria for PCOS
  \item Investigation of androgen-secreting tumors (ovarian Sertoli-Leydig cell tumors, adrenal carcinomas) --- markedly elevated testosterone >150--200 ng/dL in a woman suggests a tumor
  \item Assessment of pubertal development in males --- delayed puberty (low testosterone) vs. precocious puberty (elevated testosterone for age)
  \item Monitoring androgen deprivation therapy in prostate cancer --- target castrate testosterone <50 ng/dL, ideally <20 ng/dL
  \item Evaluation of male infertility --- testosterone is essential for spermatogenesis; low testosterone with elevated FSH suggests primary testicular failure
  \item Investigation of osteoporosis in men --- testosterone deficiency is a common cause of secondary osteoporosis in men
  \item Diagnosis of androgen insensitivity syndrome (46,XY, female phenotype) --- normal or elevated testosterone with elevated LH and no virilization
\end{itemize}

\section*{Primary vs Secondary Test}
Total testosterone is an initial test for suspected male hypogonadism and is commonly measured in the evaluation of female hyperandrogenism. In men, a diagnosis requires compatible symptoms or signs plus consistently low concentrations, confirmed with repeat early-morning testing. Free testosterone is a secondary test when total testosterone is borderline or when SHBG abnormalities are suspected (obesity, aging, diabetes, liver disease, and medications that alter SHBG) \cite{bhasin2018testosterone,wierman2014androgen}.

\section*{How This Test Is Performed}
A healthcare professional draws a blood sample from a vein into a serum separator tube (gold-top) or plain red-top tube (~3--5 mL). For suspected hypogonadism in men, collect the sample in the early morning, preferably fasting, and repeat the measurement on a separate morning before making the diagnosis. Testosterone exhibits a circadian rhythm and substantial day-to-day variation. For women being evaluated for hyperandrogenism, timing within the menstrual cycle may be documented, but clinical context and assay method are more important than applying a universal cycle-day rule. Specimen stability and handling requirements vary by laboratory.

\section*{What Preparation Is Needed}
For suspected hypogonadism in men, early-morning fasting collection is recommended; repeat a low result before diagnosis. The patient should not be acutely ill, as acute illness can transiently suppress testosterone. Follow the laboratory's instructions for biotin because interference is assay-specific. Medications that lower testosterone include opioids, glucocorticoids, GnRH agonists/antagonists, and ketoconazole. Estrogen therapy, oral contraceptives, tamoxifen, and thyroid status can alter SHBG and therefore total testosterone. Exogenous androgens suppress endogenous testosterone and LH/FSH. SHBG should be measured when free testosterone is needed.

\section*{Contraindications and Precautions}
\begin{itemize}
  \item Critical factors: testosterone MUST be measured in the morning (7--10 AM) because of the circadian rhythm. Night shift workers have a reversed rhythm and their peak testosterone occurs after daytime sleep. At least 2 morning measurements are required to confirm hypogonadism because of significant intra-individual variability. Total testosterone is profoundly affected by SHBG concentrations: conditions increasing SHBG (aging, estrogen therapy, hyperthyroidism, liver disease, anticonvulsants, HIV) increase total testosterone without affecting free testosterone; conditions decreasing SHBG (obesity, type 2 diabetes, insulin resistance, androgens, hypothyroidism, nephrotic syndrome, glucocorticoids) decrease total testosterone without affecting free testosterone. For this reason, free testosterone should be measured or calculated whenever SHBG is abnormal. Testosterone immunoassays are inaccurate at low testosterone concentrations (<150 ng/dL in men and the typical range in women/children). LC-MS/MS is the preferred measurement method per Endocrine Society and CDC guidelines, particularly for women, children, and hypogonadal men. Testosterone immunoassays cross-react with DHT (2--5\%), DHEA, androstenedione, and synthetic androgens (nandrolone, oxandrolone) to varying degrees. Acute illness, surgery, and trauma transiently suppress testosterone; testing should be delayed until recovery. Strenuous exercise can acutely elevate or suppress testosterone depending on intensity and duration. Aging is associated with a 1--2\% per year decline in total and free testosterone after age 40; age-adjusted reference ranges should be used.
\end{itemize}
\section*{Risks}
Minimal risk. Slight pain, bruising, or hematoma at venipuncture site. Rarely: vasovagal syncope. No radiation exposure.

\section*{What Happens After the Test}
Results are available according to the laboratory method. In men, a low result should be repeated in the early morning before diagnosing hypogonadism; symptoms, assay quality, and the laboratory reference interval also matter. If low testosterone is confirmed, LH and FSH help classify primary versus secondary hypogonadism. Prolactin, iron studies, or pituitary imaging are ordered selectively based on the clinical picture, not automatically for every patient. Testosterone therapy targets are formulation-specific and generally aim for a mid-normal range; treatment decisions require monitoring for adverse effects and fertility goals. In women, marked testosterone elevation or rapid virilization warrants evaluation for an androgen-secreting tumor, but no single testosterone cutoff is diagnostic across all assays. DHEA-S can help assess an adrenal source, and PCOS remains a clinical diagnosis requiring appropriate diagnostic criteria and exclusion of other causes.

\section*{Understanding Results}

\subsection*{Below Normal}
\begin{itemize}
  \item Primary hypogonadism (hypergonadotropic): low testosterone with elevated LH and FSH. Testicular failure. Causes: Klinefelter syndrome (47,XXY, most common genetic cause, 1 in 500--1000 males), cryptorchidism, orchitis (mumps), testicular torsion, chemotherapy (alkylating agents), radiation, testicular trauma, and aging (late-onset hypogonadism). Testosterone typically 100--250 ng/dL, LH >15 mIU/mL.
  \item Secondary hypogonadism (hypogonadotropic): low testosterone with inappropriately normal or low LH/FSH. Pituitary or hypothalamic failure. Causes: pituitary adenomas (prolactinoma, non-functioning), hyperprolactinemia, pituitary surgery/radiation, hemochromatosis, infiltrative diseases (sarcoidosis, histiocytosis), opioid use, anabolic steroid-induced suppression, severe obesity, critical illness, and functional hypothalamic hypogonadism from excessive exercise or weight loss. Testosterone typically <250 ng/dL, LH <5--10 mIU/mL.
  \item Androgen deprivation therapy (prostate cancer): GnRH agonists/antagonists suppress LH/FSH and testosterone to castrate levels (<50 ng/dL, target <20 ng/dL).
  \item Androgen insensitivity syndrome (46,XY): normal or elevated testosterone with elevated LH (androgen receptor defect prevents testosterone negative feedback). Complete AIS: female phenotype, testes, no uterus, 46,XY karyotype. Partial AIS: variable phenotypes.
  \item Chronic illness: systemic diseases (COPD, CKD, heart failure, HIV/AIDS, rheumatoid arthritis) suppress the hypothalamic-pituitary-gonadal axis, causing low testosterone. Approximately 25--40\% of men with type 2 diabetes have low testosterone.
  \item Medications: opioids (most common, suppress GnRH), glucocorticoids, GnRH agonists/antagonists, ketoconazole, spironolactone, and 5-alpha-reductase inhibitors.
  \item Aging (late-onset hypogonadism): progressive decline in testosterone by 1--2\% per year after age 40--50. Total testosterone declines more slowly than free testosterone because SHBG increases with age.
  \item 5-alpha-reductase deficiency (46,XY): normal testosterone but low DHT. Ambiguous genitalia at birth, virilization at puberty. Diagnosis by elevated testosterone-to-DHT ratio.
\end{itemize}

\subsection*{Above Normal}
\begin{itemize}
  \item PCOS: the most common cause of hyperandrogenism in women. Elevated total testosterone (typically 50--150 ng/dL) or free testosterone. Rotterdam criteria require 2 of 3: oligo/anovulation, clinical or biochemical hyperandrogenism, and polycystic ovarian morphology on ultrasound.
  \item Androgen-secreting ovarian tumors: Sertoli-Leydig cell tumors, hilus cell tumors, and lipid cell tumors. Rapid onset virilization, testosterone >150--200 ng/dL (often >200 ng/dL). Pelvic ultrasound or MRI detects ovarian mass. LH/FSH suppressed.
  \item Androgen-secreting adrenal tumors: adrenal adenoma or carcinoma. Testosterone and DHEA-S elevated. Adrenal CT shows adrenal mass. Typically testosterone >200 ng/dL with DHEA-S >700 mcg/dL.
  \item Congenital adrenal hyperplasia (21-hydroxylase deficiency): elevated 17-hydroxyprogesterone and androgens, low cortisol. Mild-to-moderate testosterone elevation in females with ambiguous genitalia (classical) or hirsutism/oligomenorrhea (non-classical).
  \item Exogenous testosterone/anabolic steroid use: supraphysiologic total testosterone (often >1000--2000 ng/dL) with suppressed LH/FSH and low SHBG. Erythrocytosis (hematocrit >50\%) is common.
  \item Testosterone replacement therapy: levels depend on formulation and sampling time. Results are reported in ng/dL (not ng/mL); the usual goal is a formulation-specific mid-normal range, not a universal numeric target. Interpret any high result with the dose, formulation, and time since administration.
  \item Idiopathic hirsutism: normal total and free testosterone with hirsutism. Increased 5-alpha-reductase activity in hair follicles converting normal testosterone to DHT locally.
  \item SHBG elevation: estrogen therapy, hyperthyroidism, liver disease, and HIV increase SHBG, elevating total testosterone but not free testosterone. Free testosterone measurement or calculation differentiates true hyperandrogenism.
\end{itemize}

\section*{Normal Results}
\begin{itemize}
  \item Total testosterone, adult male (19--49 years): 300--1000 ng/dL (10.4--34.7 nmol/L); age-adjusted: 40--49 years 250--900 ng/dL; 50--59 years 220--800 ng/dL; >60 years 200--700 ng/dL
  \item Total testosterone, adult female (reproductive age): 15--70 ng/dL (0.5--2.4 nmol/L); follicular phase: 15--50 ng/dL; mid-cycle: slightly higher
  \item Total testosterone, postmenopausal female: 10--40 ng/dL (0.35--1.4 nmol/L)
  \item Free testosterone, adult male: 5--21 ng/dL (0.17--0.73 nmol/L); 2\% of total
  \item Free testosterone, adult female: 0.1--0.6 ng/dL (0.003--0.02 nmol/L)
  \item Bioavailable testosterone, adult male: 100--400 ng/dL
  \item Bioavailable testosterone, adult female: 1--12 ng/dL
  \item Pre-pubertal children: total testosterone <20 ng/dL (male), <10 ng/dL (female)
  \item Castrate level (prostate cancer treatment): <50 ng/dL (GnRH agonist), ideally <20 ng/dL
  \item Common diagnostic threshold used in men: total testosterone near 300 ng/dL, confirmed on two early-morning measurements in a man with compatible symptoms or signs; use the laboratory's validated assay-specific reference interval and current guideline.
  \item Note: Reference ranges are assay- and laboratory-specific. LC-MS/MS values are lower and more accurate than immunoassay. Always use laboratory-specific ranges and note methodology.
\end{itemize}

\section*{Follow-up Tests}
\begin{itemize}
  \item LH and FSH: Classify hypogonadism as primary (elevated) or secondary (low/inappropriately normal). Essential in every hypogonadism evaluation.
  \item SHBG: Required to calculate free testosterone from total testosterone. Also explains discordant total and free testosterone levels.
  \item Free testosterone by equilibrium dialysis: Gold standard free testosterone measurement when calculated values are discrepant or when very low concentrations are expected.
  \item Prolactin: Exclude hyperprolactinemia (prolactinoma, drug-induced) as a cause of secondary hypogonadism.
  \item DHEA-S: Distinguish adrenal from ovarian hyperandrogenism in women. Elevated DHEA-S suggests adrenal source.
  \item 17-hydroxyprogesterone: Exclude non-classical congenital adrenal hyperplasia (21-hydroxylase deficiency) in hyperandrogenic women.
  \item DHT (dihydrotestosterone): Evaluate 5-alpha-reductase activity. Testosterone-to-DHT ratio >20:1 suggests 5-alpha-reductase deficiency.
  \item Karyotype: In primary testicular failure in men <50 years (Klinefelter syndrome 47,XXY) and premature ovarian insufficiency.
  \item Pituitary MRI with gadolinium: In secondary hypogonadism with low LH/FSH and low testosterone, to evaluate for pituitary adenoma, empty sella, or other sellar pathology.
  \item Iron studies (ferritin, transferrin saturation): Exclude hemochromatosis as a cause of secondary hypogonadism.
  \item Pelvic ultrasound (transvaginal): Evaluate for polycystic ovarian morphology (PCOS criteria) or ovarian androgen-secreting tumors.
  \item Semen analysis: In male infertility with abnormal testosterone/LH/FSH. Azoospermia with elevated FSH suggests non-obstructive etiology.
  \item Bone densitometry (DEXA): Screen for osteoporosis in longstanding hypogonadism.
\end{itemize}

\section*{References}
\begin{enumerate}
  \item Bhasin S, Brito JP, Cunningham GR, et al. Testosterone therapy in men with hypogonadism: an Endocrine Society clinical practice guideline. J Clin Endocrinol Metab. 2018;103(5):1715-1744.
  \item Wierman ME, Arlt W, Basson R, et al. Androgen therapy in women: an Endocrine Society clinical practice guideline. J Clin Endocrinol Metab. 2014;99(10):3489-3510.
  \item Rosner W, Auchus RJ, Azziz R, et al. Position statement: utility, limitations, and pitfalls in measuring testosterone: an Endocrine Society position statement. J Clin Endocrinol Metab. 2007;92(2):405-413.
  \item Vermeulen A, Verdonck L, Kaufman JM. A critical evaluation of simple methods for the estimation of free testosterone in serum. J Clin Endocrinol Metab. 1999;84(10):3666-3672.
  \item Teede HJ, Misso ML, Costello MF, et al. Recommendations from the international evidence-based guideline for the assessment and management of PCOS. Fertil Steril. 2018;110(3):364-379.
  \item Burtis CA, Ashwood ER, Bruns DE. Tietz Textbook of Clinical Chemistry and Molecular Diagnostics. 6th ed. Elsevier; 2023.
\end{enumerate}

% Keep the chapter citations synchronized with the maintained bibliography.
\bibliographystyle{unsrtnat}
\bibliography{references}

\clearpage
\section*{Regulatory Status and Authorization Date}
This chapter describes a test category, not a specific commercial product. FDA, CDSCO, EU IVDR, and MHRA approval or registration dates therefore cannot be assigned without the assay manufacturer, product name, instrument, and intended use. For laboratory-developed tests, an individual agency approval date may not exist. See the Preface for the interpretation of regulatory status.

\clearpage
